\subsection{COMPOSITION OF DERIVATIONS}

We suppose in this no. that case (I) of no. 2 holds, that is that A, A', A'' are three graded K-modules of type Δ and that we are given a K-bilinear mapping μ: A × A' → A'' corresponding to a graded K-linear mapping of degree 0, A ⊗K A' → A''. The graded endomorphisms f of A ⊕ A' ⊕ A'' such that f(A) ⊂ A, f(A') ⊂ A' and f(A'') ⊂ A'' form a graded subalgebra of the graded associative algebra EndgrK(A ⊕ A' ⊕ A'') (§ 3, no. 1, Example 2). In particular two ε-derivations of (A, A', A'') can be composed, but it should not be thought that the composition of two ε-derivations is an ε-derivation.

On every graded algebra B of type \(\Delta\) is defined the \(\varepsilon\) -bracket (or simply bracket when \(\varepsilon = 1\) ) of two homogeneous elements \(u, v\) , by the formula (10)

\[
[ u, v ] _ {\varepsilon} = u v - \varepsilon_ {\deg u, \deg v} v u (\text { denoted   simply   by } [ u, v ] \text { if } \varepsilon = 1).
\]

By extending this mapping by linearity, a K-bilinear mapping \((u, v) \mapsto [u, v]_{\varepsilon}\) of B × B into B is obtained. Then, for homogeneous u and v in B

\[
[ v, u ] _ {\varepsilon} = - \varepsilon_ {\deg u, \deg v} [ u, v ] _ {\varepsilon}.
\]

Applying this definition to the graded algebra \(\operatorname{Endgr}_{\mathbf{K}}(\mathbf{A} \oplus \mathbf{A}' \oplus \mathbf{A}'')\) , the \(\varepsilon\) -bracket of two graded endomorphisms is thus defined.

PROPOSITION 1. Let \(d_1, d_2\) be two \(\varepsilon\) -derivations of \((\mathbf{A}, \mathbf{A}', \mathbf{A} ' )\) of respective degrees \(\delta_1, \delta_2\) . Then the \(\varepsilon\) -bracket

\[
[ d _ {1}, d _ {2} ] _ {\varepsilon} = d _ {1} \circ d _ {2} - \varepsilon_ {\delta_ {1}, \delta_ {2}} d _ {2} \circ d _ {1}
\]

is an \(\varepsilon\) -derivation of degree \(\delta_1 + \delta_2\) . Moreover, if \(d\) is an \(\varepsilon\) -derivation of \((\mathbf{A}, \mathbf{A}', \mathbf{A}'')\) of degree \(\delta\) and if \(\varepsilon_{\delta, \delta} = -1\) , then \(d^2 = d \circ d\) is a derivation.

Suppose \(x \in A\) is homogeneous of degree \(\xi\) ; for all \(y \in A'\) ,

\[
\begin{array}{r l} d _ {1} (d _ {2} (x y)) & = ((d _ {1} d _ {2}) (x)) y + \varepsilon_ {\delta_ {1}, \delta_ {2} + \xi} (d _ {2} x) (d _ {1} y) \\ & \quad + \varepsilon_ {\delta_ {2}, \xi} (d _ {1} x) (d _ {2} y) + \varepsilon_ {\delta_ {1} + \delta_ {2}, \xi} x ((d _ {1} d _ {2}) (y)) \end{array}\tag{11}
\]

taking account of formulae (1) of no. 1. Exchanging the roles of \(d_1\) and \(d_2\) , we obtain, after simplifications again using (1) (no. 1),

\[
\begin{array}{r l} (d _ {1} d _ {2}) (x y) - \varepsilon_ {\delta_ {1}, \delta_ {2}} (d _ {2} d _ {1}) (x y) = & ((d _ {1} d _ {2}) (x)) y - \varepsilon_ {\delta_ {1}, \delta_ {2}} ((d _ {2} d _ {1}) (x)) y \\ & + \varepsilon_ {\delta_ {1} + \delta_ {2}, \xi} x ((d _ {1} d _ {2}) (y)) \\ & - \varepsilon_ {\delta_ {1}, \delta_ {2}} \varepsilon_ {\delta_ {1} + \delta_ {2}, \xi} x ((d _ {2} d _ {1}) (y)) \end{array}
\]

that is, writing \(d = [d_1, d_2]_{\varepsilon}\) and \(\delta = \delta_1 + \delta_2\) ,

\[
d (x y) = (d x) y + \varepsilon_ {\delta , \xi} x (d y)
\]

which proves that \(d\) is an \(\varepsilon\) -derivation.

On the other hand, if, in (11), we let \(d_1 = d_2 = d\) , \(\delta_1 = \delta_2 = \delta\) and \(\varepsilon_{\delta,\delta} = -1\) , we obtain, since then \(\varepsilon_{\delta,\delta + \xi} = -\varepsilon_{\delta,\xi}\) by (1),

\[
d ^ {2} (x y) = \left(d ^ {2} x\right) y + \varepsilon_ {2 5, \xi} x \left(d ^ {2} y\right)
\]

and as \(\varepsilon_{2\delta, \xi} = 1\) it is seen that \(d^2\) is a derivation.

COROLLARY. Suppose that \(\Delta = \mathbf{Z}\) . Then:

\begin{enumerate}
\def\labelenumi{(\roman{enumi})}
\item
  The square of an antiderivation is a derivation.
\item
  The bracket of two derivations is a derivation.
\item
  The bracket of an antiderivation and a derivation of even degree is an anti-derivation.
\item
  If \(d_1\) and \(d_2\) are antiderivations, \(d_1d_2 + d_2d_1\) is a derivation.
\end{enumerate}

Under the hypotheses of the beginning of this no., consider now a finite sequence \(\mathbf{D} = (d_i)_{1 \leqslant i \leqslant n}\) of pairwise permutable derivations of \((\mathbf{A}, \mathbf{A}', \mathbf{A}'')\) . For every polynomial \(\mathrm{P}(\mathbf{X}_1, \ldots, \mathbf{X}_n)\) in the algebra \(\mathbf{K}[\mathbf{X}_1, \ldots, \mathbf{X}_n]\) , the element \(\mathrm{P}(d_1, \ldots, d_n)\) of \(\operatorname{Endgr}_{\mathbf{K}}(\mathbf{A} \oplus \mathbf{A}' \oplus \mathbf{A}'')\) is then defined (§ 2, no. 9); its abbreviated notation is \(\mathrm{P}(\mathbf{D})\) .

PROPOSITION 2. With the above hypotheses and notation, consider \(2n\) indeterminates \(\mathrm{T}_1, \ldots, \mathrm{T}_n\) , \(\mathrm{T}_1', \ldots, \mathrm{T}_n'\) and for every polynomial \(\mathbf{F} \in \mathbf{K}[\mathbf{X}_1, \ldots, \mathbf{X}_n]\) write \(\mathbf{F}(\mathbf{T}) = \mathbf{F}(\mathbf{T}_1, \ldots, \mathbf{T}_n)\) , \(\mathbf{F}(\mathbf{T}') = \mathbf{F}(\mathbf{T}_1', \ldots, \mathbf{T}_n')\) and

\[
\mathrm{F} (\mathbf {T} + \mathbf {T} ^ {\prime}) = \mathrm{F} (\mathrm{T} _ {1} + \mathrm{T} _ {1} ^ {\prime}, \dots , \mathrm{T} _ {n} + \mathrm{T} _ {n} ^ {\prime}).
\]

Suppose that

\[
\mathbf {P} (\mathbf {T} + \mathbf {T} ^ {\prime}) = \sum_ {i} \mathbf {Q} _ {i} (\mathbf {T}) \mathbf {R} _ {i} (\mathbf {T} ^ {\prime})
\]

where the \(\mathbf{Q}_i\) and \(\mathbf{R}_i\) belong to \(\mathbf{K}[\mathbf{X}_1,\ldots ,\mathbf{X}_n]\) . Then, for all \(x\in \mathbf{A}\) and \(y\in \mathbf{A}'\) ,

\[
\mathbf {P} (\mathbf {D}) (x y) = \sum_ {i} \left(\mathbf {Q} _ {i} (\mathbf {D}) x\right) (\mathbf {R} _ {i} (\mathbf {D}) y).\tag{12}
\]

We introduce \(n\) other indeterminates \(T_1'', \ldots, T_n''\) and consider the polynomial algebra \(K[T_1, \ldots, T_n, T_1', \ldots, T_n', T_1'', \ldots, T_n''] = B\) ; on the other hand we consider the K-module \(M\) of bilinear mappings of \(A \times A'\) into \(A''\) ; a B-module structure is defined on \(M\) by writing, for every K-bilinear mapping \(f \in M\) and \(1 \leqslant i \leqslant n\) ,

\[
\left\{ \begin{array}{l} (\mathrm{T} _ {i} f) (a, a ^ {\prime}) = f (d _ {i} a, a ^ {\prime}) \\ (\mathrm{T} _ {i} ^ {\prime} f) (a, a ^ {\prime}) = f (a, d _ {i} a ^ {\prime}) \\ (\mathrm{T} _ {i} ^ {\prime \prime} f) (a, a ^ {\prime}) = d _ {i} (f (a, a ^ {\prime})). \end{array} \right.\tag{13}
\]

Since the \(d_{i}\) are permutable with one another, it is seen that, for every polynomial \(\mathbf{F} \in \mathbf{K}[\mathbf{X}_{1}, \ldots, \mathbf{X}_{n}]\) , \((\mathbf{F}(\mathbf{T})f)(a, a') = f(\mathbf{F}(\mathbf{D})a, a')\) ,

\[
(\mathbf {F} (\mathbf {T} ^ {\prime}) f) (a, a ^ {\prime}) = f (a, \mathbf {F} (\mathbf {D}) a ^ {\prime})
\]

and \((\mathbf{F}(\mathbf{T}^{\prime \prime})f) = \mathbf{F}(\mathbf{D})(f(a,a^{\prime}))\) . Hence, to prove (12) it suffices to show that

\[
(\mathbf {P} (\mathbf {T} ^ {\prime \prime}) - \sum_ {i} \mathbf {Q} _ {i} (\mathbf {T}) \mathbf {R} _ {i} (\mathbf {T} ^ {\prime})) \cdot \mu = 0\tag{14}
\]

or also \((\mathrm{P}(\mathbf{T}^{\prime\prime}) - \mathrm{P}(\mathbf{T} + \mathbf{T}^{\prime}))\) . \(\mu = 0\) in the B-module M. Now, the hypothesis that the \(d_{i}\) are derivations can also be expressed by saying that, for \(1 \leqslant i \leqslant n\) ,

\[
\left(\mathrm{T} _ {i} ^ {\prime \prime} - \mathrm{T} _ {i} - \mathrm{T} _ {i} ^ {\prime}\right). \mu = 0\tag{15}
\]

in the B-module M. By considering successively the polynomials

\[
\mathrm{P} \left(\mathrm{T} _ {1} ^ {\prime \prime}, \mathrm{T} _ {2} ^ {\prime \prime}, \dots , \mathrm{T} _ {n} ^ {\prime \prime}\right) - \mathrm{P} \left(\mathrm{T} _ {1} + \mathrm{T} _ {1} ^ {\prime}, \mathrm{T} _ {2} ^ {\prime \prime}, \dots , \mathrm{T} _ {n} ^ {\prime \prime}\right)
\]

\[
\mathrm{P} \left(\mathrm{T} _ {1} + \mathrm{T} _ {1} ^ {\prime}, \mathrm{T} _ {2} ^ {\prime \prime}, \dots , \mathrm{T} _ {n} ^ {\prime \prime}\right) - \mathrm{P} \left(\mathrm{T} _ {1} + \mathrm{T} _ {1} ^ {\prime}, \mathrm{T} _ {2} + \mathrm{T} _ {2} ^ {\prime}, \dots , \mathrm{T} _ {n} ^ {\prime \prime}\right)
\]

\[
\mathrm{P} \left(\mathrm{T} _ {1} + \mathrm{T} _ {1} ^ {\prime}, \dots , \mathrm{T} _ {n - 1} + \mathrm{T} _ {n - 1} ^ {\prime}, \mathrm{T} _ {n} ^ {\prime \prime}\right) - \mathrm{P} \left(\mathrm{T} _ {1} + \mathrm{T} _ {1} ^ {\prime}, \dots , \mathrm{T} _ {n - 1} + \mathrm{T} _ {n - 1} ^ {\prime}, \mathrm{T} _ {n} + \mathrm{T} _ {n} ^ {\prime}\right)
\]

it is seen that the difference \(\mathbf{P}(\mathbf{T}^{\prime \prime}) - \mathbf{P}(\mathbf{T} + \mathbf{T}^{\prime})\) can be written in the form

\[
\sum_ {i = 1} ^ {n} \left(\mathbf {T} _ {i} ^ {\prime \prime} - \mathbf {T} _ {i} - \mathbf {T} _ {i} ^ {\prime}\right) \mathrm{G} _ {i} (\mathbf {T}, \mathbf {T} ^ {\prime}, \mathbf {T} ^ {\prime \prime})
\]

where the \(G_{i}\) are elements of B. Relation (14) is therefore an immediate consequence of relations (15).

COROLLARY (Leibniz's formula). Let \(d_{i}\) ( \(1 \leqslant i \leqslant n\) ) be \(n\) derivations of \((\mathbf{A}, \mathbf{A}', \mathbf{A}'' )\) which are permutable with one another. For all \(\alpha = (\alpha_{1}, \ldots, \alpha_{n}) \in \mathbf{N}^{n}\) , we write

\[
d ^ {\alpha} = d _ {1} ^ {\alpha_ {1}} d _ {2} ^ {\alpha_ {2}} \dots d _ {n} ^ {\alpha_ {n}}.\tag{16}
\]

Then, for \(x \in \mathbf{A}\) and \(y \in \mathbf{A}'\) ,

\[
d ^ {\alpha} (x y) = \sum_ {\beta + \gamma = \alpha} ((\beta , \gamma)) d ^ {\beta} (x) d ^ {\gamma} (y)\tag{17}
\]

where we have written (in the notation introduced at the beginning of the chapter)

\[
((\beta , \gamma)) = (\beta + \gamma)! / (\beta ! \gamma!).\tag{18}
\]

This follows immediately from the multinomial formula (I, § 8, no. 2)

\[
(\mathbf {T} + \mathbf {T} ^ {\prime}) ^ {\alpha} = \sum_ {\beta + \gamma = \alpha} ((\beta , \gamma)) \mathbf {T} ^ {\beta} \mathbf {T} ^ {\prime \gamma}
\]

and Proposition 2.

